\section{Apéndice II: lista de acciones de un AI Lab ante el paradigma Agent}

Esta entrevista puede convertirse en una lista de acciones para equipos de modelos. No pretende que cada empresa la copie tal cual, sino ayudar al lector a llevar el «cambio de paradigma» a preguntas que se puedan ejecutar.

\noindent\begin{longtable}{p{0.25\textwidth}p{0.34\textwidth}p{0.32\textwidth}}
\toprule
Decisión & Pregunta que hay que responder & Práctica equivocada \\
\midrule
Marco de Agent & ¿Cómo se conecta nuestro modelo con herramientas, entornos, memoria y tareas de largo horizonte? & Tratar al Agent solo como un system prompt complejo. \\
RL infra & ¿Cómo se escalan rollout, reward, verifier y eval? & Trabajar solo con datos de preferencias fuera de línea, sin construir el ciclo cerrado con el entorno. \\
Recursos de GPU & ¿Cómo se reparten entre research, pre-train y post-train? & Seguir dejando que el preentrenamiento acapare todos los recursos estratégicos. \\
Familia de modelos & ¿Qué etapas usan modelos grandes y cuáles usan modelos pequeños o especializados? & Llamar al modelo más grande para cualquier tarea, sin distinción. \\
Costo y velocidad & ¿Cómo se optimizan en conjunto la tasa de finalización de extremo a extremo, la latencia y el precio? & Mirar solo el benchmark, sin calcular el costo de las tareas reales. \\
Estructura organizacional & ¿Entran juntos al ciclo cerrado el preentrenamiento, el posentrenamiento, el producto y los sistemas? & Dejar que un solo equipo monopolice todas las decisiones. \\
Valores & ¿Qué trabajos se sustituyen y qué valores humanos se conservan? & Optimizar solo métricas destructivas o de corto plazo. \\
\bottomrule
\end{longtable}

\begin{importantbox}{Para qué sirve la lista de acciones}
Si un equipo afirma que va a ir all in Agent, pero no puede responder las preguntas de la tabla anterior, es muy probable que solo esté siguiendo la moda. Una transformación real hacia Agent se refleja en la asignación de recursos, la infraestructura de entrenamiento, el enrutamiento entre modelos, la colaboración organizacional y la elección de valores.
\end{importantbox}

\subsection{Resumen de la sección}

El paradigma Agent no es un lema estratégico, sino un conjunto de decisiones concretas. El lector puede usar esta lista de acciones para comprobar si cualquier laboratorio de modelos, empresa emergente o equipo de producto ha entrado de verdad en la era de los Agent.
